% !TEX root = build_experiments.tex
% ---------------------------------------------------------------------------------------------
% Section 4 (Experiments), v3 2026-11-08: same four subsections as v2, content revised.
%   4.1 Experimental Setup / 4.2 Main Results / 4.3 Ablation Study / 4.4 Analysis
% Positioning (user, 2026-11-08): the contribution is segmenter-free collision constraints from scene
% geometry + VLA internals; a safety-success trade-off is acceptable and is reported, not hidden.
% Color convention:
%   black      = verified (benchmark definitions, protocol, values reported by cited papers, marked †)
%   \proj{...} = PROJECTED result, explanation, or proposed setup detail, NOT yet measured/fixed (red).
% Every number in the "our runs" blocks is projected; anchors: VLSA Tab. I (same pi0.5-LIBERO:
% CAR 17.1 % / TSR 57.8 %; AEGIS 68.9 % / 67.5 %), SafeDojo Tabs. 1-2, LIBERO-Safety Tab. 3.
% ---------------------------------------------------------------------------------------------
\section{Experiments}
\label{sec:experiments}

We evaluate \sys\ on two public safety benchmarks for VLA policies, SafeLIBERO~\cite{hu2026vlsa} and LIBERO-Safety~\cite{cui2026liberosafety}, and on a real robot, always with the same frozen $\pi_{0.5}$ policy~\cite{intelligence2025pi05} as the compared safety layers. The experiments address three questions. \textbf{Q1:} How much does \sys\ reduce the collisions of a frozen VLA while preserving its task success, and how does this balance compare with existing safety layers (Sec.~\ref{sec:exp:main})? \textbf{Q2:} Which components account for the reduction (Sec.~\ref{sec:exp:ablation})? \textbf{Q3:} Why does grounding the target from the policy's own representations work where an external segmenter does not, where does it still fail, and what does it cost (Sec.~\ref{sec:exp:analysis})?

\subsection{Experimental Setup}
\label{sec:exp:setup}

\paragraph{Benchmarks.}
SafeLIBERO~\cite{hu2026vlsa} adds obstacles to the four LIBERO suites~\cite{liu2023libero} (Spatial, Object, Goal, Long; four tasks each) at two levels: Level~I places the obstacle beside the target and Level~II on the motion path, with 50 initial states per scenario (1,600 episodes per method). The base policy is the public $\pi_{0.5}$-LIBERO checkpoint, which drives a Franka Panda through end-effector displacements at 20\,Hz under the official horizons (300 steps; 550 for Long). We also run the original LIBERO suites without obstacles to measure what the safety layer costs when nothing has to be avoided. LIBERO-Safety~\cite{cui2026liberosafety} is a separate benchmark whose physical track scores Affordance-Aware Grasping (AAG), Tabletop Spatial Avoidance (TSA), Free-Space Hand--Object Avoidance (FSHOA), and Human--Robot Interaction (HRI) at levels L0--L2, with L2 held out from training; we report L2 of AAG, TSA, and FSHOA and the full track in the supplement. The real-world study uses \proj{a Rainbow Robotics RB-Y1 with three calibrated RealSense D435 RGB-D cameras} on four tabletop tasks (clear pick-and-place, a similar object beside the target, an obstacle across the carry path, objects unseen during development), 20 trials per task and method in randomized order.

\paragraph{Baselines.}
(i)~The frozen $\pi_{0.5}$ alone. (ii)~Safety layers on the same checkpoint and evaluator: AEGIS~\cite{hu2026vlsa}, which keeps an end-effector ellipsoid away from a named obstacle with a control barrier function; multi-link CBF filtering~\cite{agarwal2026multilink}, which extends this to several links; and KNOWS~\cite{park2026knows}, which selects the target from the policy's attention but takes its extent from an external segmenter. AEGIS is run with its authors' public code; multi-link CBF and KNOWS release no code and are our re-implementations (marked~$\ast$). The obstacle is named by a VLM and an open-vocabulary detector~\cite{liu2024groundingdino} as in the original, and AEGIS is also run with the ground-truth obstacle name (GT). Two controlled baselines share our trajectory optimizer and differ only in the constraint source: \emph{generic geometry}, every observed surface an obstacle under one rule, and \emph{external segmentation}, a Grounded-SAM-2 mask~\cite{liu2024groundingdino,ravi2024sam2} prompted with the instruction's noun, lifted through the same depth, and combined with our contact policy. (iii)~Values reported by the authors of VLSA and KNOWS for the same $\pi_{0.5}$-LIBERO checkpoint under their own evaluator runs (marked~$\dagger$), which let the reader check our reproduction; safety-trained policies on OpenVLA-OFT~\cite{kim2025openvlaoft,zhang2025safevla,zhu2026wmpo,tang2026safedojo} use a different base policy and are reported in the supplement. Reported values are never ranked against our runs.

\paragraph{Metrics and implementation.}
We follow each benchmark: SafeLIBERO reports task success rate (TSR), collision avoidance rate (CAR; no active obstacle displaced by more than 1\,mm), and safe success rate (SSR; success and no collision in the same episode); LIBERO-Safety counts a success (SR) only without a safety violation and reports the collision rate (CR). On the real robot, unsafe contacts are scored from \proj{a wrist force--torque sensor and two external cameras reviewed by annotators blind to the method}, never from the method's own estimate. \sys\ reads the noun attention from \proj{layer~8, head~6} of $\pi_{0.5}$ and the identity features from \proj{prefix layer~8}; these and all thresholds were fixed on \proj{the first five initial states of each scenario} and every reported number comes from \proj{the remaining 45 (1,440 episodes per method)}. The scene distance field uses cuRobo~\cite{sundaralingam2023curobo} and the chunk is refined by SQP with OSQP~\cite{stellato2020osqp} through the robot's Jacobian, since LIBERO policies output end-effector displacements (Sec.~3.5). All arms share initial states and policy seeds, every failure, timeout, and hold stays in the denominator, and paired differences carry 95\% cluster-bootstrap intervals over scenarios; further details are in the supplement.

\subsection{Main Results}
\label{sec:exp:main}

\begin{table*}[t]
\centering
\caption{\textbf{Results on LIBERO and SafeLIBERO (\%).} All rows use the frozen $\pi_{0.5}$-LIBERO checkpoint. Upper block ($\dagger$): values reported by the authors of VLSA~\cite{hu2026vlsa} (Tab.~I, suite average, levels pooled, no SSR) and KNOWS~\cite{park2026knows} (Tab.~1; we average their four suite rows and report CAR as $100-$CR, where their CR counts any contact with a non-target obstacle) under their own evaluator runs; Naive is their ellipsoid filter on ground-truth segmentation. Lower block: our runs with the official evaluator on the same episodes for every row; $\ast$ our re-implementation (no public code), GT a privileged obstacle name. LIBERO is the original suite without obstacles. Reported values are not ranked against our runs; best of the lower block in bold. \proj{Red values are projected targets, not measurements.}}
\label{tab:safelibero}
\footnotesize
\setlength{\tabcolsep}{4pt}
\begin{tabular}{llc|ccc|ccc}
\toprule
 & & LIBERO & \multicolumn{3}{c|}{SafeLIBERO Level I} & \multicolumn{3}{c}{SafeLIBERO Level II}\\
Method & Obstacle input & SR$\uparrow$ & TSR$\uparrow$ & SSR$\uparrow$ & CAR$\uparrow$ & TSR$\uparrow$ & SSR$\uparrow$ & CAR$\uparrow$\\
\midrule
\multicolumn{9}{l}{\emph{Reported by the authors, same checkpoint}$^\dagger$}\\
$\pi_{0.5}$ alone~\cite{hu2026vlsa} & -- & -- & \multicolumn{6}{c}{TSR 57.8 \quad SSR -- \quad CAR 17.1 \quad (levels pooled)}\\
AEGIS~\cite{hu2026vlsa} & VLM + detector & -- & \multicolumn{6}{c}{TSR 67.5 \quad SSR -- \quad CAR 68.9 \quad (levels pooled)}\\
$\pi_{0.5}$ alone~\cite{park2026knows} & -- & -- & 55.0 & 11.5 & 12.4 & 57.8 & 20.0 & 23.0\\
Naive (GT segmentation)~\cite{park2026knows} & privileged & -- & 73.0 & 63.1 & 82.2 & 67.8 & 55.6 & 76.1\\
KNOWS~\cite{park2026knows} & attention + segmenter & -- & 67.5 & 59.9 & 76.4 & 71.1 & 53.8 & 68.8\\
\midrule
\multicolumn{9}{l}{\emph{Our runs, frozen $\pi_{0.5}$-LIBERO, same evaluator}}\\
$\pi_{0.5}$ alone & -- & \proj{\textbf{96.8}} & \proj{57.8} & \proj{13.5} & \proj{18.6} & \proj{59.4} & \proj{19.0} & \proj{25.2}\\
+ generic geometry & depth & \proj{93.9} & \proj{42.3} & \proj{39.1} & \proj{\textbf{88.4}} & \proj{44.0} & \proj{41.2} & \proj{\textbf{90.1}}\\
+ external segmentation & depth + mask & \proj{95.0} & \proj{60.0} & \proj{46.5} & \proj{71.0} & \proj{61.5} & \proj{48.0} & \proj{74.0}\\
AEGIS (official code) & VLM + detector & \proj{96.5} & \proj{66.5} & \proj{49.8} & \proj{69.0} & \proj{67.0} & \proj{51.0} & \proj{72.0}\\
AEGIS (official code, GT) & privileged & \proj{96.5} & \proj{\textbf{67.5}} & \proj{52.5} & \proj{73.5} & \proj{\textbf{68.0}} & \proj{53.5} & \proj{75.5}\\
Multi-link CBF$^\ast$~\cite{agarwal2026multilink} & VLM + detector & \proj{96.1} & \proj{65.0} & \proj{50.0} & \proj{72.5} & \proj{66.0} & \proj{51.5} & \proj{75.0}\\
KNOWS$^\ast$~\cite{park2026knows} & attention + segmenter & \proj{95.8} & \proj{66.5} & \proj{55.0} & \proj{76.0} & \proj{69.0} & \proj{53.0} & \proj{70.0}\\
\textbf{\sys\ (ours)} & depth + VLA internals & \proj{95.4} & \proj{66.0} & \proj{\textbf{61.5}} & \proj{87.5} & \proj{\textbf{68.0}} & \proj{\textbf{63.0}} & \proj{89.0}\\
\bottomrule
\end{tabular}
\end{table*}

\begin{table}[t]
\centering
\caption{\textbf{LIBERO-Safety, physical track, L2 (\%).} SR counts episodes completed without a safety violation; CR is pooled over the three suites. Upper block ($\dagger$): reported in~\cite{cui2026liberosafety}, mean\,$\pm$\,std over three training seeds. Lower block: our runs with the single public $\pi_{0.5}$ weight and the official evaluator. \proj{Red values are projected.}}
\label{tab:liberosafety}
\scriptsize
\setlength{\tabcolsep}{2pt}
\begin{tabular}{lcccc}
\toprule
Method & AAG & TSA & FSHOA & CR$\downarrow$\\
\midrule
\multicolumn{5}{l}{\emph{Reported}$^\dagger$}\\
OpenVLA~\cite{kim2024openvla} & 4.0$\pm$0.0 & 12.7$\pm$0.9 & 0.0$\pm$0.0 & --\\
OpenVLA-OFT~\cite{kim2025openvlaoft} & 1.3$\pm$0.9 & 40.0$\pm$1.6 & 42.7$\pm$2.5 & --\\
$\pi_0$~\cite{black2024pi0} & 11.3$\pm$2.5 & 0.7$\pm$0.9 & 0.7$\pm$0.9 & --\\
$\pi_{0.5}$~\cite{intelligence2025pi05} & 35.3$\pm$2.5 & 56.7$\pm$3.8 & 51.3$\pm$6.1 & --\\
GR00T N1.6~\cite{nvidia2025gr00t} & 19.3$\pm$3.4 & 52.7$\pm$2.5 & 50.0$\pm$1.6 & --\\
\midrule
\multicolumn{5}{l}{\emph{Our runs, same evaluator}}\\
$\pi_{0.5}$ alone & \proj{32.7} & \proj{54.0} & \proj{49.3} & \proj{26.0}\\
+ generic geometry & \proj{19.3} & \proj{57.3} & \proj{56.0} & \proj{8.7}\\
+ external segmentation & \proj{29.3} & \proj{56.7} & \proj{53.3} & \proj{14.0}\\
\textbf{\sys\ (ours)} & \proj{\textbf{37.3}} & \proj{\textbf{62.0}} & \proj{\textbf{60.7}} & \proj{\textbf{6.7}}\\
\bottomrule
\end{tabular}
\end{table}

\textbf{SafeLIBERO.} Tab.~\ref{tab:safelibero} reports four-suite averages. \proj{Our run of the frozen policy alone completes 57.8\% of Level~I episodes and avoids the obstacle in 18.6\%, between the 17.1\% and 12.4\% CAR reported for the same checkpoint by~\cite{hu2026vlsa} and~\cite{park2026knows}: nothing in the policy represents the added object, so it executes the demonstrated motion through it. Avoiding the collision also completes more tasks, because a pushed obstacle often topples the target or blocks the placement; every safety layer, reported or run by us, raises TSR over the policy alone, and \sys\ raises it from 57.8\% to 66.0\%. The decomposition shows where the groups differ: the policy's TSR\,$-$\,SSR gap of 44.3 points is success achieved \emph{through} a collision, which \sys\ removes (CAR 87.5\%, SSR 61.5\%); its own CAR\,$-$\,SSR gap of 26.0 points is collision-free episodes that still fail, 79\% of them in a hold that times out near the obstacle, and on obstacle-free LIBERO the same holds cost 1.4 points (95.4\% vs.\ 96.8\%). Among the safety layers on the same episodes, KNOWS$^\ast$, the strongest without privileged input, reaches a CAR of 76.0\% and an SSR of 55.0\%, close to its authors' 76.4\% and 59.9\%; \sys\ collides in 12.5\% of episodes against 24.0\%, about half as often, exceeds even AEGIS with the ground-truth obstacle name by 14.0 CAR points, and keeps TSR within 1.5 points of every layer. The gap has two sources, which Sec.~\ref{sec:exp:ablation} isolates. The CBF layers protect an ellipsoid of the end-effector, or of a few links, against the one obstacle that was named, so collisions of the other hand links and of the attached object with unnamed clutter remain (61\% of the residual collisions of AEGIS); \sys\ keeps the spheres of both hands and of the attached object clear of every observed surface. And on Level~I the obstacle touches the target, where an ellipsoid or an external mask either swallows the target and blocks the grasp or frees the obstacle with it; the depth-continuous surface keeps the two apart, which is why \sys\ comes within two SSR points of the Naive filter of~\cite{park2026knows} although that filter is given ground-truth segmentation. Generic geometry shows the opposite failure, CAR 88.4\% at TSR 42.3\%, because the target itself is forbidden. On Level~II \sys\ reaches an SSR of 63.0\%, whereas KNOWS loses SSR from Level~I to Level~II both in its authors' runs (59.9\% to 53.8\%) and in ours: a segmenter or a learned safety signal is tuned to the obstacle placement it was developed on, our constraints are read from the observed scene at every request.}

\textbf{LIBERO-Safety.} Tab.~\ref{tab:liberosafety} tests the same mechanism against a different safety taxonomy. \proj{AAG rewards contact with the target and penalizes contact with its surroundings, exactly the allowed-versus-forbidden split of our contact policy: \sys\ raises SR from 32.7\% to 37.3\%, whereas generic geometry drops to 19.3\% because the grasp is forbidden. TSA and FSHOA are the generic-obstacle case; \sys\ reaches 62.0\% and 60.7\% against 54.0\% and 49.3\%, and the pooled collision rate falls from 26.0\% to 6.7\%. Our $\pi_{0.5}$ row is below the reported 35.3 / 56.7 / 51.3 because the public weight is a single training seed.}

\begin{figure}[t]
\centering
\fbox{\parbox[c][3.6cm][c]{0.96\columnwidth}{\centering\proj{Figure to prepare: 4 task columns (clear / similar neighbor / carry-path obstacle / unseen objects) $\times$ 2 rows ($\pi_{0.5}$ alone, \sys), each a third-person frame at the moment of contact or avoidance with a wrist-view inset showing the selected surface (orange) and contact links (gold); below, a bar chart of task and safe success per task, 20 trials each.}}}
\caption{\textbf{Real-world results.} The same frozen $\pi_{0.5}$ with and without \sys\ on \proj{the RB-Y1}, from calibrated RGB-D only; 20 trials per task and method. Unsafe contacts are scored from external sensing. \proj{Values projected.}}
\label{fig:real}
\end{figure}

\textbf{Real-world evaluation.} Fig.~\ref{fig:real} reports 20 trials per task and method. \proj{The frozen policy completes 71.3\% of trials but makes an unsafe contact in 45.0\%, most often by dragging the attached object across the obstacle or by pressing the neighbor while grasping. \sys\ completes 68.8\% with unsafe contacts in 8.8\%; generic geometry completes 55.0\% with 11.3\%. On objects unseen during development, \sys\ keeps 62.5\% task success and 10.0\% unsafe contacts (policy alone: 65.0\% and 42.5\%).}

\subsection{Ablation Study}
\label{sec:exp:ablation}

\begin{table}[t]
\centering
\caption{\textbf{Ablation on SafeLIBERO Level~I (\%).} One component of \sys\ is replaced by its nearest alternative; everything else, including the optimizer, is unchanged. $\Delta$SSR is the paired difference to the full method with a 95\% cluster-bootstrap interval. \proj{Red values are projected.}}
\label{tab:ablation}
\scriptsize
\setlength{\tabcolsep}{2pt}
\begin{tabular}{lcccc}
\toprule
Configuration & TSR & SSR & CAR & $\Delta$SSR [95\% CI]\\
\midrule
\textbf{\sys\ (full)} & \proj{66.0} & \proj{\textbf{61.5}} & \proj{87.5} & --\\
Seeds: attention $\rightarrow$ nearest object & \proj{62.0} & \proj{54.8} & \proj{81.5} & \proj{$-$6.7 [$-$9.1, $-$4.2]}\\
Extent: graph $\rightarrow$ radius growth & \proj{65.0} & \proj{54.3} & \proj{77.8} & \proj{$-$7.2 [$-$9.8, $-$4.5]}\\
Extent: graph $\rightarrow$ external mask & \proj{60.0} & \proj{51.2} & \proj{72.0} & \proj{$-$10.3 [$-$13.0, $-$7.4]}\\
Identity: memory off & \proj{64.5} & \proj{57.4} & \proj{83.5} & \proj{$-$4.1 [$-$6.2, $-$1.9]}\\
Policy: target always obstacle & \proj{36.0} & \proj{34.5} & \proj{\textbf{93.0}} & \proj{$-$27.0 [$-$30.5, $-$23.2]}\\
Policy: target freed, all links/stages & \proj{\textbf{66.5}} & \proj{55.5} & \proj{77.0} & \proj{$-$6.0 [$-$8.4, $-$3.5]}\\
Policy: no object attachment & \proj{65.5} & \proj{54.6} & \proj{78.0} & \proj{$-$6.9 [$-$9.3, $-$4.4]}\\
\bottomrule
\end{tabular}
\end{table}

Tab.~\ref{tab:ablation} replaces one component at a time on SafeLIBERO Level~I; every variant keeps the optimizer, so the differences isolate the representation and the contact policy.

\textbf{Target identity.} \proj{Replacing the attention seeds by the object nearest to the hand costs 6.7 SSR points, almost entirely in the Object suite, where a distractor is closer than the target in about a third of the initial states: the rule then frees the wrong object and forbids the right one. Turning the object memory off costs 4.1 points, concentrated in episodes where the hand occludes the target at the moment of grasp and the per-request selection jumps to a neighbor.}

\textbf{Target extent.} \proj{Growing the seeds by radius instead of along depth-continuous surfaces costs 7.2 points and lifting an external mask costs 10.3, in both cases through CAR rather than TSR: the freed region takes the touching obstacle with it, and the optimizer no longer sees it. This is the largest effect among the perception components and the one an external segmenter cannot repair, because it segments in the image, where touching objects share a boundary.}

\textbf{Contact policy.} \proj{Keeping the target an obstacle in every stage gives the highest CAR (93.0\%) but cuts TSR by 30 points, since the grasp is forbidden; freeing it for every link in every stage restores TSR but drops CAR to 77.0\%, since the wrist and the attached object may then enter the target region and its surroundings. Leaving the grasped object unattached costs 6.9 points, all in the transport phase, where the object rather than the hand hits the obstacle. All three terms are needed to reach the full method.}

\subsection{Analysis}
\label{sec:exp:analysis}

\begin{table}[t]
\centering
\caption{\textbf{Target grounding on held-out scenes.} Soft 3D IoU, precision, and recall of the predicted observed surface against simulator ground truth, target-identification accuracy (ID), and identity switches per sequence (Sw.) in the occlusion sequences; averaged frame\,$\rightarrow$\,episode\,$\rightarrow$\,task. The segmenter is prompted with the noun of the instruction unless noted; YOLOE and a privileged oracle are in the supplement. \proj{Red values are projected.}}
\label{tab:grounding}
\scriptsize
\setlength{\tabcolsep}{3pt}
\begin{tabular}{lccccc}
\toprule
Method & IoU & Prec. & Rec. & ID & Sw.$\downarrow$\\
\midrule
Geometry, nearest-to-hand rule & \proj{0.44} & \proj{0.61} & \proj{0.58} & \proj{0.63} & \proj{--}\\
Grounded-SAM-2~\cite{ravi2024sam2} & \proj{0.48} & \proj{0.55} & \proj{0.74} & \proj{0.81} & \proj{4.2}\\
Grounded-SAM-2, canonical name & \proj{0.61} & \proj{0.66} & \proj{0.80} & \proj{0.90} & \proj{3.1}\\
\midrule
\sys\ w/o memory & \proj{0.69} & \proj{0.88} & \proj{0.76} & \proj{0.90} & \proj{2.8}\\
\textbf{\sys\ (ours)} & \proj{\textbf{0.72}} & \proj{\textbf{0.90}} & \proj{0.78} & \proj{\textbf{0.96}} & \proj{\textbf{0.4}}\\
\bottomrule
\end{tabular}
\end{table}

\textbf{Why policy-internal grounding beats an external segmenter.} Tab.~\ref{tab:grounding} scores the target surface on held-out diagnostic scenes with a single target, similar objects, touching clutter, and partial occlusion. \proj{\sys\ reaches an IoU of 0.72 at a precision of 0.90, against 0.48 for Grounded-SAM-2 prompted with the instruction and 0.61 when it is given the object's canonical name. The segmenter's recall matches ours, but its precision is 0.55: when a neighbor touches the target the mask spans both, and when the robot occludes the target the mask follows the robot (Fig.~1, panel~d). The grounding is conditioned on the instruction the policy executes: changing only the instruction on an identical scene and robot state moves the seed, the surface, and the resulting constraint to the newly named object in \proj{n/N} scene pairs. Our precision comes from the depth-continuous graph, which stops at every visible depth break, and from the rendered robot mask. The segmenter also depends on the prompt, which costs 0.13 IoU between the canonical name and the instruction's noun, whereas the attention is conditioned on the instruction the policy actually executes; this is why KNOWS$^\ast$, which selects by attention but segments externally, inherits the leakage of Tab.~\ref{tab:safelibero}. Identity is a temporal question: the memory reduces identity switches from 2.8 to 0.4 per sequence and raises identification accuracy to 0.96 while leaving single-frame IoU almost unchanged.}

\textbf{Where it still fails.} \proj{Classifying every collision and every collision-free failure of \sys\ on Level~I (supplement, Tab.~S9) splits them into two groups. Perception-side causes account for 71\% of the 12.5\% collisions: surfaces that were never observed (46\%), a target identity that drifted from the scene under occlusion or clutter (17\%), and a hand surface outside the collision spheres (8\%). Action-side causes account for the remaining collisions and for 93\% of the collision-free failures: a corridor narrower than the margin ends in a hold (79\%), and a deflected chunk disturbs a grasp or placement pose that the optimizer cannot regenerate (14\%). The first group is a limit of how well the policy's representation is synchronized with the observed geometry, the second of refining a chunk after it is generated; Sec.~5 takes up both.}

\textbf{Generalization and run time.} \proj{The only learned component is the frozen policy, so the shift \sys\ is exposed to is perceptual: on held-out objects, layouts, and occlusions withheld from every threshold choice, its SSR drops by 1.8 points against 7.5 for KNOWS$^\ast$ and 9.1 for the external-segmentation baseline, whose detectors miss the unseen objects. The two fixed choices are not fragile: the three attention heads ranked next on the development set change held-out IoU by at most \proj{0.02}, and a sphere margin of $\pm$5\,mm moves SSR by \proj{$\pm$1.1} points (supplement, Tab.~S14). Run time is measured at one boundary, from observation to an executable chunk, including the policy's own inference (\proj{about 200}\,ms on one H200). \sys\ answers in \proj{490}\,ms (p50) / \proj{620}\,ms (p95), against \proj{205}\,ms for the policy with AEGIS, whose CBF-QP adds under a millisecond per step~\cite{hu2026vlsa}, and \proj{560}\,ms for the external-segmentation baseline; of our cost the attention read-out, taken from the forward pass that produces the actions, is \proj{6}\,ms, the surface graph \proj{95}\,ms, the two distance fields \proj{80}\,ms, and the refinement \proj{110}\,ms (supplement, Tab.~S11). The precise surface and the link-specific fields thus cost one to two orders of magnitude more than an ellipsoid filter. We report this as a cost, not as real-time capability, and return to it in Sec.~5.}
